\documentclass[12pt]{article}
\usepackage[utf8]{inputenc}
\usepackage[T1]{fontenc}
\usepackage{hyperref}
\title{Uma' Lung Language Resources for K-12 Curriculum Development}
\author{Compiled by Claude Sonnet 5 (claude-sonnet-5) for Andr\'as Kornai}
\date{}
\begin{document}
\maketitle

\section{Introduction}
Uma' Lung (ISO 639-3 ulu; Glottocode umal1238) is an endangered Austronesian (Kenyah) language of Borneo, Indonesia, with about 3,000 speakers. Searches found very little: no dictionary, grammar, primer or school materials were located. Note that Uma' Lasan (xky) is a separate Sarawak variety and should not be confused with it.

\section{Primary Educational Resources}
No primers, textbooks or school programmes for Uma' Lung were found in search results.

\section{Linguistic and Reference Materials}
\begin{itemize}
\item Glottolog languoid page for Uma' Lung: \url{https://glottolog.org/resource/languoid/id/umal1238.bigmap.html}
\item Endangered Languages Project entry (a profile noting that no dedicated dictionary or grammar exists): \url{https://endangeredlanguages.com/node/40474}
\item Endangered Languages Project bibliography for the language: \url{https://endangeredlanguages.com/node/40474/bibliography}
\item Linguisten.de blog post ``Endangered Languages Challenge: Uma' Lung'': \url{https://blog.linguisten.de/post/186883236459/endangered-languages-challenge-uma-lung}
\item WALS entry (code keu): \url{https://wals.info/languoid/lect/wals_code_keu}
\item Wikipedia page for the related Uma' Lasan language, useful for disambiguation: \url{https://en.wikipedia.org/wiki/Uma%CA%BC_Lasan_language}
\end{itemize}
Search summaries also mention a published story book (Soriente 2006, Atma Jaya University Press, Jakarta) and scholarly articles by Blust (2007) and Smith (2015) in \emph{Oceanic Linguistics}; no URLs for these appeared in the results, so they are not listed as resources here.

\section{Religious and Community Resources}
None found.

\section{Digital and Audio Resources}
A documentation profile reports some elicitation and story recordings exist, but no accessible audio collection was found in the results.

\section{Curriculum Development Context}
Materials for a native-language curriculum essentially do not yet exist in findable form. A curriculum would need to begin with community documentation: a wordlist, an orthography, and recorded stories. The Soriente story collection and Blust's phonology are leads to pursue through library channels.
\end{document}